\documentclass[fleqn]{article}
\usepackage[utf8]{inputenc}
\usepackage[T2A]{fontenc}
\usepackage[russian]{babel}
\usepackage{graphicx}
\usepackage{amsmath}
\renewcommand{\arraystretch}{1.5}
\renewcommand{\baselinestretch}{1.75}
\usepackage[margin=3cm]{geometry}

\usepackage{amssymb}
\usepackage{caption}
\setlength{\mathindent}{0pt}

\begin{document}

\title{Задачка}
\author{Цуканников Кирилл Андреевич}
\maketitle

\newpage
Задача: четверо игроков подбрасывают монету, на которой с одинаковой вероятностью 0,5 выпадает «орел» или «решка». Выигрывает тот участник, у которого впервые выпадает «орел». Определить вероятность выигрыша каждого из участников. \\
Решение: \\
Пусть:\\
 А - победа 1го участника \\
 B - победа 2го участника \\
 C - победа 3го участника \\
 D - победа 4го участника \\
В задаче явно не написано, после проигрыша 4 участника идут ли они по кругу или наступает, пожтому я рассмотрю 2 случая: после проигрыша 4го игра заканчивается участники играют по кругу до 1 победы. \\
1) После 4го игра заканчивается \\
P(A) = 0.5 \\ (Выпал или не выпал орел)
P(B) = 0.5 * 0.5 \\(у 1го не выпал и у 2го либо выпал, либо нет. Следующие аналогично)
P(C) = 0.5 * 0.5 * 0.5 \\
P(D) = 0.5 * 0.5 * 0.5 * 0.5 \\
Ответ: P(A) = 0.5 \\
P(B) = 0.25 \\
P(C) = 0.0625 \\
P(D) = 0.003125 \\
2) После 4го игра продолжается с 1го \\
В таком случае речь идет об геометрическом распределении, так как вероятность постоянна (0.5) и игра происходит до первого благоприятного события \\
В теории игра может происходить бесконечно, поэтому рассмотрю сумму до $+\infty$ \\
1й игрок может победить на 1, 5, 9, ... шагах. \\
Пусть  p - вероятность победы, q = 1 - p- проигрышь \\
Хотя обычно геометрическое распределние начинается с 1, а не с 0, тут я использую с 1, чтобы можно было 1 круг рассмотреть. \\
Тогда: \\
$ P(A) = \displaystyle\sum_{m=0}^{\infty}pq^{(4m+1) - 1} =\displaystyle\sum_{m=0}^{\infty}0.5 * 0.5^{4m} = 0.5 \displaystyle\sum_{m=0}^{\infty}(\frac{1}{16})^m $  - геометрическая прогрессия\\
$S_n = \dfrac{b_1(q^n) - 1}{q-1}$ \\
$S_{\infty} = \dfrac{1((\frac{1}{16})^{\infty} - 1) }{\frac{1}{16}-1} = \frac{16}{15}$ \\ 
Ответ: \\
$ P(A)  = 0.5 *  \frac{16}{15} = \frac{8}{15} $ \\
$ P(B)  = 0.5 * \frac{8}{15} = \frac{4}{15} $ \\
$ P(C)  = 0.5 * \frac{4}{15} = \frac{2}{15} $ \\
$ P(D)  = 0.5 * \frac{2}{15} = \frac{1}{15} $ \\



\end{document}